\chapter{パーミッションとユーザ管理}
\label{chap:permission}

% この章で学ぶこと
\begin{goalbox}
  \begin{itemize}
    \item \cmd{ls -l} で表示されるパーミッション（権限）の読み方
    \item \cmd{chmod} と \cmd{chown} を使った権限と所有者の変更
    \item \cmd{sudo} を安全に使うための考え方
    \item センサやマイコンなどのデバイスを使うための権限の設定
  \end{itemize}
\end{goalbox}

\ref{sec:linux-internals-users}節で学んだように、Linux のすべてのファイルには所有者がいて、誰がそのファイルを読んだり書いたりできるかが決められています。
この「誰が何をしてよいか」の設定を\textbf{パーミッション}（権限）といいます。
「\cmd{Permission denied}（権限がありません）」というエラーは、ロボット開発で最もよく目にするエラーの 1 つです。
この章では、パーミッションのしくみを理解し、このエラーに自分で対処できるようになりましょう。

\section{パーミッションの読み方}
\label{sec:permission-read}

\subsection{\cmd{ls -l} の表示}

ファイルのパーミッションは、\cmd{ls -l}（\ref{sec:files-navigation}節）で確認できます。
この章では、\ref{chap:files}章で作った \cmd{\textasciitilde/practice} の中に、練習用のディレクトリ \cmd{perm} を作って試します。

\begin{terminal}[練習用のファイルを作って、パーミッションを確認する]
$ mkdir ~/practice/perm
$ cd ~/practice/perm
$ mkdir scripts
$ touch notes.txt start_robot.sh
$ ls -l
total 4
-rw-rw-r-- 1 user user    0 Sep 30 10:01 notes.txt
drwxrwxr-x 2 user user 4096 Sep 30 10:00 scripts
-rw-rw-r-- 1 user user    0 Sep 30 10:02 start_robot.sh
\end{terminal}

左端の 10 文字がパーミッションで、3 列目と 4 列目がそれぞれ所有者とグループです（日時やファイルの大きさは、環境によって変わります）。
10 文字は、図\ref{fig:permission-string}のように区切って読みます。

\begin{figure}[tb]
  \centering
  % 図：パーミッションの表記の読み方（chapters/commandline/permission.tex）
\begin{tikzpicture}[ch/.style={font=\Large\ttfamily, inner sep=0pt, outer sep=0pt, anchor=base west, minimum width=5mm}]
  \foreach \c [count=\i from 0] in {-,r,w,x,r,-,x,r,-,-}
    \node[ch] (c\i) at (\i*0.5,0) {\c};
  % 区切りの色
  \begin{scope}[on background layer]
    \fill[black!10] ([yshift=-1mm]c0.south west) rectangle ([yshift=5mm]c0.south east);
    \fill[blue!15] ([yshift=-1mm]c1.south west) rectangle ([yshift=5mm]c3.south east);
    \fill[green!15] ([yshift=-1mm]c4.south west) rectangle ([yshift=5mm]c6.south east);
    \fill[orange!15] ([yshift=-1mm]c7.south west) rectangle ([yshift=5mm]c9.south east);
  \end{scope}
  \node[figlabel, below=3mm of c0] {種類};
  \node[figlabel, below=3mm of c2, align=center] {所有者\\（user）};
  \node[figlabel, below=3mm of c5, align=center] {グループ\\（group）};
  \node[figlabel, below=3mm of c8, align=center] {その他\\（others）};
  % 右側の説明
  \node[figlabel, anchor=west, align=left] at (5.6,-0.2)
    {\texttt{-}：通常のファイル　\texttt{d}：ディレクトリ\\
     \texttt{r}：読み取り（read）\\
     \texttt{w}：書き込み（write）\\
     \texttt{x}：実行（execute）\\
     \texttt{-}：その権限がない};
\end{tikzpicture}

  \caption{パーミッションの表記の読み方（\cmd{-rwxr-xr--} の場合）}
  \label{fig:permission-string}
\end{figure}

\begin{enumerate}
  \item 1 文字目はファイルの種類です。\cmd{-} は通常のファイル、\cmd{d} はディレクトリ、\cmd{l} はシンボリックリンク、\cmd{c} はデバイスファイル（\ref{sec:linux-internals-device}節）を表します。
  \item 2〜4 文字目は、\textbf{所有者}（user）に対する権限です。
  \item 5〜7 文字目は、ファイルに設定された\textbf{グループ}（group）のメンバーに対する権限です。
  \item 8〜10 文字目は、それ以外のすべてのユーザ（\textbf{その他}、others）に対する権限です。
\end{enumerate}

それぞれの 3 文字は、左から順に「読み取り（\cmd{r}）」「書き込み（\cmd{w}）」「実行（\cmd{x}）」ができるかどうかを表し、できない場合は \cmd{-} になります。
今作った \cmd{notes.txt} の \cmd{-rw-rw-r--} は、所有者とグループのメンバーが読み書きでき、その他のユーザは読み取りだけができる、という意味です。
図\ref{fig:permission-string}の \cmd{-rwxr-xr--} なら、次の意味になります（次の節で、\cmd{start\_robot.sh} をこの権限に変えてみます）。

\begin{itemize}
  \item 所有者の \cmd{user} は、読み取り・書き込み・実行のすべてができる（\cmd{rwx}）
  \item グループ \cmd{user} のメンバーは、読み取りと実行ができる（\cmd{r-x}）
  \item その他のユーザは、読み取りだけができる（\cmd{r--}）
\end{itemize}

\subsection{ファイルとディレクトリでの意味の違い}

\cmd{r}, \cmd{w}, \cmd{x} の意味は、ファイルとディレクトリで少し異なります。

\begin{tablebox}[ファイルとディレクトリでの権限の意味][tab:permission-rwx]
  \begin{tabular}{lp{0.35\linewidth}p{0.35\linewidth}}
    \toprule
    権限 & ファイルの場合 & ディレクトリの場合 \\
    \midrule
    \cmd{r} & 中身を読める & 中のファイルの一覧を見られる \\
    \cmd{w} & 中身を書き換えられる & 中にファイルを作ったり、削除したりできる \\
    \cmd{x} & プログラムとして実行できる & \cmd{cd} で中に入れる \\
    \bottomrule
  \end{tabular}
\end{tablebox}

シェルスクリプト（\ref{chap:shellscript}章）や Python のスクリプト（\ref{sec:python-run}節）を \cmd{./start\_robot.sh} のように直接実行するには、ファイルに \cmd{x} の権限が必要です。
権限がないと、次のようなエラーになります。

\begin{terminal}[実行権限がないファイルを実行した場合]
$ ./notes.txt
bash: ./notes.txt: Permission denied
\end{terminal}

\section{chmod, chown の使い方}
\label{sec:permission-chmod}

\subsection{\cmd{chmod}：権限を変える}

ファイルの権限を変えるには、\cmd{chmod}（change mode）を使います。
権限の指定の仕方には、記号で書く方法と数字で書く方法があります。

\textbf{記号で指定する方法}では、「誰に（\cmd{u}：所有者、\cmd{g}：グループ、\cmd{o}：その他、\cmd{a}：全員）」「追加するか削除するか（\cmd{+} / \cmd{-}）」「どの権限を（\cmd{r}, \cmd{w}, \cmd{x}）」を組み合わせて書きます。

\begin{terminal}[記号で権限を変える]
$ chmod u+x start_robot.sh      # 所有者に実行権限を追加
$ chmod go-w notes.txt          # グループとその他から書き込み権限を削除
$ chmod +x start_robot.sh       # 誰にも指定しないと、全員に追加
\end{terminal}

最もよく使うのは、スクリプトに実行権限を付ける \cmd{chmod +x} です。

\textbf{数字で指定する方法}では、\cmd{r} を 4、\cmd{w} を 2、\cmd{x} を 1 として、所有者・グループ・その他のそれぞれについて足し算した数字を 3 桁並べて書きます。

\begin{tablebox}[\cmd{chmod} の数字による指定の例][tab:permission-numeric]
  \begin{tabular}{llll}
    \toprule
    数字 & 表記 & 意味 & よく使う場面 \\
    \midrule
    \cmd{755} & \cmd{rwxr-xr-x} & 所有者はすべて、ほかは読み取りと実行 & スクリプト、ディレクトリ \\
    \cmd{644} & \cmd{rw-r--r--} & 所有者は読み書き、ほかは読み取りだけ & 通常のファイル \\
    \cmd{700} & \cmd{rwx------} & 所有者だけがすべてできる & 自分専用のディレクトリ \\
    \cmd{600} & \cmd{rw-------} & 所有者だけが読み書きできる & 秘密鍵など \\
    \bottomrule
  \end{tabular}
\end{tablebox}

\begin{terminal}[数字で権限を変える]
$ chmod 754 start_robot.sh
$ ls -l start_robot.sh
-rwxr-xr-- 1 user user 0 Sep 30 10:02 start_robot.sh
$ chmod 755 start_robot.sh
$ ls -l start_robot.sh
-rwxr-xr-x 1 user user 0 Sep 30 10:02 start_robot.sh
\end{terminal}

たとえば \cmd{755} は、所有者が $4+2+1=7$、グループが $4+1=5$、その他が $4+1=5$ という意味です。
\cmd{754} なら、その他が $4$（読み取りだけ）になり、図\ref{fig:permission-string}と同じ \cmd{-rwxr-xr--} になります。

\begin{caution}[\cmd{chmod 777} は使わない]
  「\cmd{Permission denied} が出たら \cmd{chmod 777} すればよい」という情報をインターネットで見かけることがあります。
  \cmd{777} は、誰でも読み書き・実行ができるという設定で、ほかのユーザや悪意のあるプログラムに書き換えられる危険があります。
  エラーが出たら、まずは \cmd{ls -l} で「誰に」「どの権限が」足りないのかを確かめて、必要な権限だけを追加しましょう。
\end{caution}

\subsection{\cmd{chown}：所有者を変える}

ファイルの所有者やグループを変えるには、\cmd{chown}（change owner）を使います。
所有者を変えられるのは root だけなので、\cmd{sudo} を付けて実行します。
\cmd{所有者:グループ} のように書くと、両方を同時に変えられます。
ディレクトリの中身もまとめて変えるときは \cmd{-R} を付けます。

うっかり \cmd{sudo} を付けて作ったファイルは、所有者が root になり、自分では編集できなくなります（\ref{sec:linux-internals-users}節）。
そのようなときに、\cmd{chown} で所有者を自分に戻します。
\cmd{sudo} を付けてファイルを作り、所有者を戻してみましょう。

\begin{terminal}[所有者を変える]
$ sudo touch root_file.txt
$ ls -l root_file.txt
-rw-r--r-- 1 root root 0 Sep 30 10:05 root_file.txt
$ sudo chown user:user root_file.txt
$ ls -l root_file.txt
-rw-r--r-- 1 user user 0 Sep 30 10:05 root_file.txt
\end{terminal}

\cmd{user:user} の部分は、自分のユーザ名に置き換えてください（\cmd{\$USER:\$USER} と書くこともできます。\cmd{\$USER} については\ref{chap:env}章で説明します）。
ディレクトリを中身ごと変えるときは、\cmd{sudo chown -R user:user ディレクトリ名} のように \cmd{-R} を付けます。

\section{sudo と安全な権限の扱い方}
\label{sec:permission-sudo}

\subsection{\cmd{sudo} のしくみ}

\ref{sec:linux-internals-users}節で説明したように、\cmd{sudo} は、後ろに書いたコマンドを root の権限で実行するためのコマンドです。
\cmd{sudo} を使えるのは \cmd{sudo} グループのメンバーだけで、どのユーザがどのコマンドを実行できるかは \cmd{/etc/sudoers} というファイルで決められています。
一度パスワードを入力すると、しばらく（標準では 15 分）はパスワードを聞かれずに \cmd{sudo} を使えます。

\subsection{\cmd{sudo} を使うべきとき・使うべきでないとき}

\cmd{sudo} が必要になるのは、主にシステム全体に関わる操作をするときです。

\begin{itemize}
  \item ソフトウェアのインストールや更新（\cmd{sudo apt install} など、\ref{chap:package}章）
  \item \cmd{/etc} の下の設定ファイルの編集
  \item ほかのユーザが所有するファイルの操作
\end{itemize}

逆に、自分のホームディレクトリの中での作業には、\cmd{sudo} は必要ありません。
ROS 2 のワークスペースのビルド（\ref{chap:workspace}章）や、自分のプログラムの実行、Python のパッケージのインストール（\ref{chap:python_env}章）に \cmd{sudo} を付けると、root が所有するファイルができてしまい、後でトラブルの原因になります。

\begin{point}[\cmd{Permission denied} が出たときの考え方]
  \begin{enumerate}
    \item エラーの出たファイルやディレクトリを \cmd{ls -l} で確認し、所有者・グループ・権限を見る
    \item 自分のファイルなのに権限が足りないなら、\cmd{chmod} で必要な権限を追加する
    \item 所有者が root になってしまっているなら、\cmd{sudo chown} で自分に戻す
    \item システムの設定を変える操作なら、\cmd{sudo} を付けて実行する
  \end{enumerate}
  「とりあえず \cmd{sudo} を付ける」のは、根本的な解決にならないことが多いので避けましょう。
\end{point}

\subsection{\cmd{sudo} とリダイレクト}

\cmd{sudo} を使っても、うまくいかない書き方があります。
たとえば、root しか書き込めないファイルに、\cmd{>}（\ref{sec:text-pipe-redirect}節）で文字を書き込もうとすると、\cmd{sudo} を付けてもエラーになります。

\begin{terminal}[sudo を付けても書き込めない例]
$ sudo echo "hello" > /etc/robot.conf
bash: /etc/robot.conf: Permission denied
\end{terminal}

これは、\cmd{sudo} が root の権限で実行するのは \cmd{echo} だけで、ファイルへの書き込み（\cmd{>}）は、一般ユーザの権限のシェルが行っているためです。
このような場合は、\cmd{tee} コマンドを \cmd{sudo} で実行して書き込みます。

\begin{terminal}[tee を使って root の権限で書き込む]
$ echo "hello" | sudo tee /etc/robot.conf
hello
\end{terminal}

\cmd{tee} は、受け取った文字を画面に表示しながらファイルにも書き込むコマンドです（\ref{sec:text-pipe-pipe}節）。
Docker のリポジトリを登録するとき（\ref{sec:docker-basic}節）にも、この書き方が使われています。

\section{デバイスへのアクセス権}
\label{sec:permission-device}

\subsection{デバイスファイルの権限}

USB でつないだセンサやマイコンを使おうとして、\cmd{Permission denied} というエラーが出ることがあります。
デバイスファイル（\ref{sec:linux-internals-device}節）のパーミッションを見てみましょう。
次は、USB シリアルの機器（マイコンボードなど）をつないだときの例です。
機器をつないでいなければ \cmd{/dev/ttyUSB0} はないので、同じ結果にはなりません。

\begin{termexample}[デバイスファイルの権限を確認する]
$ ls -l /dev/ttyUSB0
crw-rw---- 1 root dialout 188, 0 Sep 30 10:15 /dev/ttyUSB0
\end{termexample}

1 文字目の \cmd{c} は、デバイスファイル（キャラクタデバイス）であることを表しています。
所有者は root、グループは \cmd{dialout} で、所有者とグループのメンバーだけが読み書きできる（\cmd{rw-rw----}）ことがわかります。
つまり、このデバイスを使うには、自分が \cmd{dialout} グループのメンバーである必要があります。

\subsection{グループにユーザを追加する}

自分が所属しているグループは、\cmd{groups} コマンドで確認できます。

\begin{terminal}[所属しているグループを確認する]
$ groups
user adm cdrom sudo dip plugdev users lpadmin
\end{terminal}

\cmd{dialout} が含まれていなければ、\cmd{usermod} コマンドで自分を \cmd{dialout} グループに追加します。
\cmd{-aG} は、「今のグループに加えて（append）、このグループ（Group）にも所属させる」という意味です。
\cmd{\$USER} は、自分のユーザ名が入った環境変数です（\ref{chap:env}章）。

\begin{terminal}[dialout グループに自分を追加する]
$ sudo usermod -aG dialout $USER
\end{terminal}

\textbf{グループの変更は、一度ログアウトしてログインし直すまで反映されません。}
再ログインしてから \cmd{groups} を実行し、\cmd{dialout} が含まれていることを確かめましょう。

\begin{caution}[\cmd{-a} を忘れない]
  \cmd{usermod -G dialout} のように \cmd{-a} を付け忘れると、今まで所属していたグループ（\cmd{sudo} など）から外れて、\cmd{dialout} だけに所属することになってしまいます。
  \cmd{sudo} グループから外れると \cmd{sudo} が使えなくなるので、必ず \cmd{-aG} と書きましょう。
\end{caution}

ロボット開発では、\cmd{dialout} のほかに、次のようなグループがよく出てきます。

\begin{tablebox}[ロボット開発でよく出てくるグループ][tab:permission-groups]
  \begin{tabular}{lp{0.6\linewidth}}
    \toprule
    グループ & 使えるようになるもの \\
    \midrule
    \cmd{dialout} & USB シリアルの機器（\cmd{/dev/ttyUSB0}, \cmd{/dev/ttyACM0}） \\
    \cmd{video} & カメラ（\cmd{/dev/video0}） \\
    \cmd{input} & ゲームパッドなどの入力機器（\cmd{/dev/input/}） \\
    \cmd{docker} & \cmd{sudo} なしでの \cmd{docker} コマンド（\ref{chap:docker}章） \\
    \bottomrule
  \end{tabular}
\end{tablebox}

\begin{column}{udev でデバイスの設定を固定する}
  \cmd{sudo chmod 666 /dev/ttyUSB0} のようにデバイスファイルの権限を直接変えても、機器を挿し直したり再起動したりすると元に戻ってしまいます。
  また、デバイスファイルの番号は、つないだ順番で変わってしまいます（\ref{sec:linux-internals-device}節）。
  実際のロボットでは、\textbf{udev} というしくみを使って、これらを解決することがよくあります。
  udev は、機器がつながったときに、デバイスファイルの名前や権限を決めるしくみです。
  \cmd{/etc/udev/rules.d/} の下に\textbf{ルール}を書いたファイルを置いておくと、たとえば「この LiDAR がつながったら、いつでも \cmd{/dev/lidar} という名前を付け、誰でも使える権限にする」といった設定ができます。
  センサのメーカーが、udev のルールファイルを用意してくれていることも多いので、機器の説明書を確認してみましょう。
\end{column}

\section*{この章のまとめ}

\begin{itemize}
  \item パーミッションは、所有者・グループ・その他のそれぞれに対して、読み取り（\cmd{r}）・書き込み（\cmd{w}）・実行（\cmd{x}）ができるかを表します。\cmd{ls -l} で確認できます。
  \item \cmd{chmod} で権限を、\cmd{sudo chown} で所有者を変えます。スクリプトには \cmd{chmod +x} で実行権限を付けます。\cmd{chmod 777} は避けましょう。
  \item \cmd{sudo} はシステム全体に関わる操作のときだけ使います。\cmd{Permission denied} が出たら、まず \cmd{ls -l} で原因を確かめます。
  \item USB のセンサやマイコンを使うには、\cmd{dialout} グループに自分を追加します（反映には再ログインが必要）。
  \item udev のルールを使うと、デバイスに固定の名前と権限を設定できます。
\end{itemize}
